\section{Implications and Future Work}\label{sec:conclusion}
In this paper, we have introduced a notion of local elasticity for neural networks. This notion enables us to develop a new clustering algorithm, and its effectiveness on the MNIST and CIFAR-10 datasets provide evidence in support of the local elasticity phenomenon in neural networks. While having shown the local elasticity in both synthetic and real-world datasets, we acknowledge that a mathematical foundation of this notion is yet to be developed. Specifically, how to rigorous formulate this notion for neural networks? A good solution to this question would involve the definition of a meaningful similarity measure and, presumably, would need to reconcile the notion with possible situations where the dependence between the prediction change and the similarity is not necessarily
monotonic. Next, can we prove that this phenomenon occurs under some geometric structures of the dataset? Notably, recent evidence suggests that the structure of the training data has a profound impact on the performance of neural networks \citep{Modellingtheinfluence}. Moreover, how does local elasticity depend on network architectures, activation functions, and optimization strategies?

Broadly speaking, local elasticity implies that neural networks can be plausibly thought of as a local method. We say a method is local if it seeks to fit a data point only using observations within a \textit{window} of the data point. Important examples include the $k$-nearest neighbors algorithm, kernel smoothing, local polynomial regression \citep{fan2018local}, and locally linear embedding \citep{roweis2000nonlinear}. An exciting research direction is to formally relate neural networks to local methods. As opposed to the aforementioned classic local methods, however, neural networks seem to be capable of choosing the right \textit{bandwidth} (window size) by adapting the data structure. It would be of great interest to show whether or not this adaptivity is a consequence of the elasticity of neural networks.

%An exciting research direction is to formally relate neural networks to local methods, and to understand how neural networks are often capable of adapting to the right \textit{bandwidth} (window size) in complex real-life problems. 
In closing, we provide further implications of this notion by seeking to interpret various aspects of neural networks via local elasticity. Our discussion lacks rigor and hence much future investigation is needed.

{\bf Memorization.} Neural networks are empirically observed to be capable of fitting even random labels perfectly \citep{zhang2016understanding}, with provable guarantees under certain conditions \citep{allen2019convergence, du2019gradient2,oymak2019towards}. Intuitively, local elasticity leads neural nets to progressively fit the examples in a vicinity of the input via each SGD update, while remaining fitting (most) labels that have been learned previously. A promising direction for future work is to relate local elasticity to memorization in a more concrete fashion.

{\bf Stability and generalization.} \cite{bousquet2002stability} demonstrate that the uniform stability of an algorithm implies generalization on a test dataset. As noted by \cite{kuzborskij2018data}, due to its distribution-free and worst-case nature, uniform stability can lead to very loose bounds on generalization. The local elasticity of neural networks, however, suggests that the replace of one training point by another imposes limited perturbations on the predictions of most points and thus the loss might be more stable than expected. In this regard, it is possible to introduce a new notion of stability that relies on the metric of the input space for better generalization bounds. 

{\bf Data normalization and batch normalization.} The two normalization techniques are commonly used to improve the performance of neural networks \citep{gonzalez2002digital,ioffe2015batch}. The local elasticity viewpoint appears to imply that these techniques allow for a more solid relationship between the relative prediction change and a certain distance between feature vectors. Precisely, writing the lifting map $\varphi(\bx) \approx \frac{\partial f(\bx, \bw)}{\partial \bw}$, Section~\ref{subsec:kernel} reveals that the kernelized similarity approximately satisfies $2\Sk(\mbf{x}, \mbf{x}') \approx \|\varphi(\mbf{x})\|^2 + \|\varphi(\mbf{x}')\|^2 - d(\mbf{x}, \mbf{x}')^2$, where $d(\mbf{x}, \mbf{x}') = \|\varphi(\mbf{x}) - \varphi(\mbf{x}')\|$. To obtain a negative correlation between the similarity and the distance, therefore, one possibility is to have the norms of $\varphi(\mbf{x})$ and $\varphi(\mbf{x}')$ about equal to a constant. This might be made possible by employing the normalization techniques. See some empirical results in Appendix \ref{sec:disc-some-terms}. A venue for future investigation is to consolidate this implication on normalization techniques.
%Finally, we provide a self-contained justification of the behaviors of neural nets on the meomorization, generalization, robustness, large-scale data, and the trick of data normalization and batch normalization. The simple interpretation is shown in the following paragraphs, and further verification is postponed as our future work.
%%, demonstrate the efficiency of the proposed clustering algorithm. It further supports our theory about the local elasticity of neural nets. 
%%propose a new geometric interpretation, the local elasticity, to explain the effectiveness of neural nets. We first show that neural nets with nonlinear activation functions have the property of local elasticity by extensive simulations. 
%% Moreover, can we develop a concept similar to bandwith in nonparametric statistics that can be used to quantitatively characterize how the prediction change decreases as the geodesic distance? 
%%It is our hope that more interested researchers could further develop the notion of local elasticity toward unveiling the effectiveness of neural networks. 

%%also known as data interpolation. Because of the local elasticity, neural nets will mainly update the points in the local neighborhood of the updated point. In this way, the predictions of points are not easily influenced by the SGD updates on other points, so neural nets can easily memorize the training data. 
%%Memorization is necessary for generalization \citep{feldman2019does}.

%{\bf Implicit regularization.} This aspect of neural networks refers to the observation that SGD updates is prone to a path that regularizes the weights in a certain fashion without imposing explicit regularization. \weijie{to cite} To provide insights into implicit regularization using local elasticity, note that Figure~\ref{fig:simulations} shows that the relative similarity seems to be a smooth function of the geodesic distance of two feature vectors, hence the name elasticity. Thus, the decision boundary or, more broadly, the prediction function is likely to be smooth in the input space, an outcome of implicit regularization.
%% This intimate connection between generalization is extended to models that are trained using SGD \citep{hardt2016train}. \weijie{cite more}

%%\weijie{to discuss KNN more. doesn’t explain why certain architectures generalize better than other architectures. different architecutres have different levels of local elasticity.}
%the Because the SGD update of neural nets is like a elastic string, it will not only influence the prediction of the updated point but also those of the points in the local neighborhood. In this way, neural nets can generalize well in the local neighborhood. 

%%{\bf Robustness.} Because of the local elasticity of neural nets, the parameters updated on noisy data points cannot affect a lot on other correct data points. In this way, neural nets are robust on noises.

%%{\bf Large-Scale Data.} There are many manifolds composed of the data points in the large-scale data. We can update the neural nets with any order, because the updates of points in one manifold have a small effect on the predictions of points in other manifolds because of the local elasticity. As for the linear models, we should sample randomly from all manifolds to guarantee the performance, because the updates of points in one manifold easily affect the predictions of points in other manifolds with the linear models.  However, it is almost impossible to sample randomly from all manifolds because we don't know the latent manifolds. In general, it is easier for neural nets to learn the large scale data with many manifolds.

%\hangfeng{Dataset and control dataset is confusing, because you can use dataset in general case.}

%%This explanation can support the fact that data normalization \citep{gonzalez2002digital,ioffe2015batch} and batch normalization \citep{ioffe2015batch} can improve the performance of neural nets. We further compare the clustering algorithm w/o data normalization and find that data normalization is also important for our clustering algorithm to have a good performance. More details can be found in Appendix \ref{subsec:data-normalization}. 

%%As discussed in Section \ref{subsec:kernel}, the local elasticity of neural nets is based on the assumption that $S_{ker}(\mbf{x}, \mbf{x}') = \frac{\norm{\varphi(\mbf{x})}^2 + \norm{\varphi(\mbf{x}')}^2 - d(\mbf{x}, \mbf{x}')}{2}$ has a negative relation with the geodesic distance $d(x, x')$ between $x$ and $x'$ in the latent manifold. It actually requires that the norm of $\varphi(x)$ and $\varphi(x')$ is close to the constant. Data normalization and batch normalization can help keep the norm of latent features of neural nets close to the constant, because the norm of the input or hidden vectors can directly affect the norm of the latent features. For example, the norm of $x$ will easily affect the norm of $\frac{\partial f(w, x)}{\partial w}$ in the two-layer ReLU activated neural nets. This explanation can support the fact that data normalization \citep{gonzalez2002digital,ioffe2015batch} and batch normalization \citep{ioffe2015batch} can improve the performance of neural nets. We further compare the clustering algorithm w/o data normalization and find that data normalization is also important for our clustering algorithm to have a good performance. More details can be found in Appendix \ref{subsec:data-normalization}. 



%%Future work involves various directions. We list a few here. First, we plan to apply the proposed clustering algorithms in the multiple classes setting. Second, it is interesting to develop a semi-supervised learning algorithm based on the local elasticity to make use of the annotated resources. We may be able to make use of the large scale data for pre-training and then fine-tune it on the specific data like BERT \citep{devlin2019bert}. Finally, we want to provide rigorous theoretical guarantees for the local elasticity of neural nets. Actually, we can use the spectral clustering to further improve the method but we postpone it as our future work. 

%%\weijie{To mathematically formulate the definition of local elasticity. mention bandwith. batch size. batch size is effectively 1 here. stability, and implicit regularization.} 


%%% Local Variables:
%%% mode: latex
%%% TeX-master: "paper"
%%% End:
